\section{Introduction}
\label{sec:intro}

Robot demonstrations are expensive, while egocentric human recordings and other auxiliary embodiments are cheap and abundant. Mixed training can exploit this asymmetry, but auxiliary data can also produce negative transfer. Choosing a source fraction requires understanding the target risk across mixing weights: a locally favorable weight need not be globally best, and finite samples can change which candidate is preferred.

We study a dual-domain model with Gaussian covariates, a shared spectral basis, task-vector mismatch $\delta$, and domain-specific noise. Standard full-ridge regression is the main estimator; diagonalized ridge is retained only as an interpretable special case. The model is deliberately simple so that the risk geometry and the global optimum can be characterized rather than inferred from a local derivative or a coarse grid. Closed-loop robot policies are outside the scope of this paper.

\paragraph{Two coordinates already produce multiple valleys.}
Our starting point is the population risk as a function of the mixing weight. One coordinate cannot produce multiple valleys, whereas two suffice (Proposition~\ref{prop:multimodal}). A certified two-dimensional example and a one-parameter family exhibit local minima separated by a local maximum, valley creation and annihilation, and two distinct global minimizers. This makes the location and value of every candidate relevant to selecting the mixture. Zero mismatch supplies a contrasting case: log-convexity gives a complete uniqueness and flatness classification.

\paragraph{Algebraic global certification.}
The population derivative can be cleared of positive denominators to obtain a polynomial stationary equation of degree at most $3p-2$. For fixed parameters, complete real-root isolation followed by comparison of all stationary and boundary candidates certifies the global optimum without assuming convexity or unimodality. This supplies a direct way to check when the nondegenerate perturbation theorem applies and when a different global-shape analysis is required.

\paragraph{Finite-sample optimal selection.}
A fixed-dimension expansion of the actual full-ridge risk and its first two derivatives connects these population cases to finite-sample optima (Theorem~\ref{thm:t1c}). At a unique nondegenerate interior population optimum, it gives a first-order weight correction and optimized risk and gain (Theorem~\ref{thm:optfull}). With exactly two nondegenerate interior population minimizers, the smaller $R_1$ branch wins for sufficiently large samples. When matched spectra and zero mismatch make $R_0$ flat, $R_1=A(1-w)^2+Bw^2$ determines the limiting weight and first-order risk and gain if $A+B>0$. If $A=B=0$, the risk is exactly zero for every weight. A first-order tie between distinct nontrivial branches is unresolved at this order.

\paragraph{Evidence and scope.}
Exact root isolation and interval comparisons certify population geometry. Conditional Monte Carlo checks the actual-risk expansion, interior weight corrections, and flat-case gains. The reported comparison at two fixed population weights tests their finite-sample risk ordering; it does not validate persistence or selection of finite-sample local-minimum branches. A general finite-sample shape classification, parameter-estimation guarantees, and robot-policy validation remain open.

\paragraph{Contributions.}
We analyze global risk geometry and its consequences for choosing a mixing weight:
(i)~a minimum-dimension multimodality result and a certified two-coordinate family with valley bifurcations and two global minima;
(ii)~polynomial global certification for fixed parameters, together with a zero-mismatch uniqueness and flatness classification;
(iii)~a uniform fixed-dimension risk and derivative expansion yielding nondegenerate weight corrections, first-order selection between two population minima, and flat-population optimal selection.
Proofs, exact parameters, and numerical protocols are in the supplementary material.
